% ==============================================================================
% FFGFT: Dokument 246 (Deutsch)
% Titel: RSG und FFGFT — Ein struktureller Vergleich
% Autor: Johann Pascher
% Datum: Mai 2026
% ==============================================================================

\section*{Zusammenfassung}

Recursive Survival Geometry (RSG, Austin 2024--2026) und die Fundamentale
Fraktale Geometrische Feldtheorie (FFGFT/T0, Pascher 2024--2026) behandeln
überlappende Fragen — welche physikalische Struktur persistiert und warum —
von strukturell verschiedenen Ausgangspunkten aus. Dieses Dokument vergleicht
die beiden Frameworks entlang vier Dimensionen: was jedes Framework als
Primitiv nimmt, was es ableitet, wo die beiden Ansätze komplementär sind,
und wo FFGFT tiefer geht im Sinne numerischer Vorhersagen, die RSG nicht
liefern kann. Das Ergebnis: die beiden Frameworks sind asymmetrisch
komplementär. FFGFT liefert die Skalenwerte, die RSGs Filter als Eingabe
benötigt aber nicht ableitet; RSG formuliert das Selektionsprinzip, das
FFGFT instantiiert aber nicht explizit als allgemeinen Operator formuliert.
Keines ist ein Overlay oder eine Brückenschicht des anderen.

% ------------------------------------------------------------------------------
\section{Was RSG ist}
% ------------------------------------------------------------------------------

Recursive Survival Geometry ist Peter Austins Framework für beobachtbare
physikalische Struktur als differentielle Persistenz von Histories unter
einem rekursiven überlebensgewichteten Filter. Die Zentralgleichung lautet:

\begin{equation}
  \frac{dS_i}{dt} = -\Gamma_i \cdot W_i \cdot S_i
\end{equation}

wobei $S_i$ das Überlebensgewicht von History $i$ ist, $\Gamma_i$ die
Verlustrate und $W_i$ der Überlebensfaktor. Was wir als Materie, Strahlung
oder Struktur beobachten, sind die Histories, die unter diesem Filter
persistieren.

RSGs ontologisches Primitiv ist die \emph{History} — ein Pfad durch den
Konfigurationsraum — und der \emph{Überlebensfilter} — der Operator, der
Persistenz gewichtet. Der Filter ist allgemein: RSG legt sich nicht auf
einen bestimmten Lagrangian, einen bestimmten Feldinhalt oder eine bestimmte
Menge von Kopplungskonstanten fest.

\subsection{Was RSG ableitet}

RSG leitet die strukturellen Bedingungen ab, unter denen eine History
persistieren kann: sie muss die wiederholte Anwendung des Filters überleben,
Abschluss unter dem rekursiven Prozess erhalten und Gewicht $W_i$ über der
Schwelle akkumulieren. Der Formalismus liefert eine allgemeine Beschreibung
warum manche Konfigurationen persistieren und andere nicht.

\subsection{Was RSG nicht ableitet}

RSG leitet keine Werte physikalischer Konstanten ab. Die Verlustraten
$\Gamma_i$, die Überlebensgewichte $W_i$ und die Schwellenbedingungen sind
Framework-Parameter. RSG erklärt nicht warum die Elektronenmasse
$0{,}511\,\text{MeV}$ beträgt, warum die Feinstrukturkonstante
$\alpha \approx 1/137$ ist, oder warum die CMB-Temperatur $2{,}725\,\text{K}$
beträgt. Das sind Eingaben für RSGs Filter, keine Ausgaben.

% ------------------------------------------------------------------------------
\section{Was FFGFT ist}
% ------------------------------------------------------------------------------

Die Fundamentale Fraktale Geometrische Feldtheorie (T0/Zeit-Masse-Dualität)
nimmt den 4-dimensionalen Torus $T^4 = (\mathbb{R}/2\pi\mathbb{Z})^4$ mit
Wicklungsstruktur $(n_\theta, n_\phi, n_3, n_4)$ als ontologisches Primitiv.
Die Grundrelation ist $\tilde{T} \cdot m = 1$ (Zeit-Masse-Dualität). Ein
einziger dimensionsloser Parameter $\xi = 4/30000$ wird aus drei gemessenen
Massenverhältnissen über den eindeutigen Wicklungsexponenten $\kappa = 7$
bestimmt.

Der T0-Lagrangian lautet:

\begin{align}
  \mathcal{L}_{\text{T0}} &= -\tfrac{1}{4}F_{\mu\nu}F^{\mu\nu}
    + \bar{\psi}_\ell(i\gamma^\mu D_\mu - m_\ell)\psi_\ell \nonumber\\
    &\quad + \tfrac{1}{2}(\partial_\mu \Delta m)(\partial^\mu \Delta m)
    - \tfrac{1}{2}m_T^2 \Delta m^2
    + \xi \cdot m_\ell \cdot \bar{\psi}_\ell \psi_\ell \cdot \Delta m
\end{align}

wobei $\Delta m(x,t)$ das dynamische Massenfeld und $\xi \cdot m_\ell$ der
Vertex-Kopplungsfaktor ist.

\subsection{Was FFGFT ableitet}

Aus $\xi = 4/30000$ allein:

\begin{itemize}
  \item Lepton- und Quark-Massenverhältnisse via $\kappa=7$-Rekursion
        (Fehler $<1\%$, kein Fitting)
  \item Feinstrukturkonstante $\alpha$ (Dok.~009/137)
  \item CMB-Temperatur $T_\text{CMB}$ (Abweichung $\Delta = 0{,}0002\%$, Dok.~025)
  \item Kosmologische Konstante $\Lambda = 1{,}34\times10^{-122}$
        (Planck-Einheiten, Dok.~182)
  \item Hubble-Konstante $H_0 = 66{,}2\,\text{km/s/Mpc}$ (Dok.~182)
  \item Geometrischer Rotverschiebungsmechanismus durch Pfadverlängerung
        (FEM-Simulation, Dok.~026)
\end{itemize}

\subsection{Was FFGFT nicht explizit formuliert}

FFGFT hat kein explizites Selektionsprinzip vom RSG-Typ. Es leitet die
Skalenstruktur des Existierenden ab, formuliert aber nicht warum diese
Strukturen als beobachtbare Histories persistieren statt im Vakuum
absorbiert zu werden. Die Persistenzfrage — warum Lepton-artige Histories
überleben und andere nicht — wird implizit durch die Wicklungsstruktur von
$T^4$ beantwortet, aber nicht durch einen dedizierten Überlebensfilterformalismus.

% ------------------------------------------------------------------------------
\section{Der strukturelle Vergleich}
% ------------------------------------------------------------------------------

\renewcommand{\arraystretch}{1.5}
{\small
\begin{longtable}{p{3.5cm}p{5.0cm}p{5.0cm}}
\toprule
\textbf{Dimension} & \textbf{RSG (Austin)} & \textbf{FFGFT (Pascher)} \\
\midrule
\endfirsthead
\toprule
\textbf{Dimension} & \textbf{RSG (Austin)} & \textbf{FFGFT (Pascher)} \\
\midrule
\endhead
\midrule\multicolumn{3}{r}{\footnotesize\itshape Fortsetzung\ldots}\\\endfoot
\bottomrule\endlastfoot

\textbf{RA4 Primitiv} &
History + Überlebensfilter. Kein festgelegter Feldinhalt oder Lagrangian. &
$T^4$-Wicklungsgeometrie. Festgelegter Lagrangian: $\mathcal{L}_{\text{T0}}$. \\

\midrule

\textbf{RA3 Zulässigkeit} &
Histories müssen rekursive Filteranwendung über Schwelle überstehen.
Allgemeine Bedingung, keine fixen Kopplungskonstanten. &
Nur mit $\xi$ konsistente Wicklungskonfigurationen zulässig.
$D_f = 3-\xi$, Z3$^3$-Bedingung fixiert drei Generationen. \\

\midrule

\textbf{RA2 Darstellung} &
Gewichteter History-Raum mit Überlebensmaß $S_i$. Kein fixer Hilbertraum. &
$L^2(T^4)$; Fourier-Moden mit Wicklungszahlen als natürliche Basis;
QM-Operatoralgebra als Projektion hergeleitet (Dok.~230/231). \\

\midrule

\textbf{RA1.5 Operat.} &
$\Gamma_i$, $W_i$ als Framework-Parameter — nicht aus fixierter Eingabe
abgeleitet. Skala ohne Zusatzeingabe undefiniert. &
Alle SM-Konstanten aus $\xi$ allein: Massen, $\alpha$, $T_\text{CMB}$,
$\Lambda$, $H_0$. Kein Fitting. \\

\midrule

\textbf{RA1 Vorhersagen} &
Strukturelle Persistenzbedingungen. Keine spezifischen numerischen
Vorhersagen für Teilchenmassen oder kosmologische Konstanten. &
Konkrete Zahlenwerte für alle SM-Konstanten und kosmologischen
Parameter. Explizite Falsifizierungskriterien. \\

\midrule

\textbf{Rotverschiebung} &
Im RSG-Kern nicht explizit adressiert. &
Geometrische Pfadverlängerung im $\xi$-Feld; FEM-Simulation bestätigt
Frequenzunabhängigkeit. Von innen ununterscheidbar von metrischer
Expansion (Dok.~026). \\

\midrule

\textbf{Selektionsprinzip} &
Explizit: Überlebensfilter $\Gamma_i W_i$ definiert welche Histories
persistieren. &
Implizit: Wicklungszahl-Zulässigkeit. Kein dedizierter Persistenzoperator
separat vom Lagrangian formuliert. \\

\midrule

\textbf{Skalenableitung} &
Nicht vorhanden: $\Gamma_i$, $W_i$ benötigen externe Eingabe für
spezifische Werte. &
Vollständig: $\xi$ aus drei Massenverhältnissen; alle Skalen folgen
ohne weitere Eingabe. \\

\end{longtable}
}

% ------------------------------------------------------------------------------
\section{Ist RSG durch FFGFT erklärbar?}
% ------------------------------------------------------------------------------

Teilweise. FFGFT liefert die Skalenwerte, die RSGs Filter als Eingabe
benötigt. Wenn $\Gamma_i = \xi \cdot m_i$ (Kopplung proportional zur Masse,
wie im $\mathcal{L}_{\text{T0}}$-Vertex), dann sind RSGs Filter-Parameter
durch FFGFTs Kopplungsstruktur bestimmt. In dieser Lesart ist RSGs
Überlebensfilter die phänomenologische Beschreibung dessen, was FFGFTs
Lagrangian auf der Ebene der beobachtbaren History-Selektion erzeugt.

Was FFGFT nicht liefert, ist RSGs allgemeine Formulierung des
Selektionsprinzips als framework-unabhängiger Operator. RSGs Überlebensfilter
gilt für jede History unabhängig davon, ob die zugrundeliegende Dynamik
von FFGFT, von Standard-QFT oder von einer anderen Feldtheorie kommt.
Diese Allgemeinheit ist RSGs Beitrag; FFGFTs spezifische Kopplungsstruktur
ist eine Instantiierung davon.

% ------------------------------------------------------------------------------
\section{Ist FFGFT durch RSG erklärbar?}
% ------------------------------------------------------------------------------

Nein. RSG leitet keine Zahlenwerte für Teilchenmassen, Kopplungskonstanten
oder kosmologische Parameter ab. Die Elektronenmasse $m_e = 0{,}511\,\text{MeV}$,
das Verhältnis $m_p/m_e = 1836{,}15$, die Feinstrukturkonstante
$\alpha \approx 1/137$ — keines davon folgt aus RSGs Überlebensfilter
ohne externe Spezifikation von $\Gamma_i$ und $W_i$. RSGs Formalismus ist
mit jedem Wert dieser Konstanten konsistent; er sagt sie nicht vorher.

FFGFT hingegen fixiert $\xi$ aus drei gemessenen Verhältnissen und leitet
alle anderen ab. Das ist ein stärkerer Anspruch, kein schwächerer. FFGFT
ist kein rechnerisches Overlay von RSG — es ist eine konkurrierende und
stärker eingeschränkte Antwort auf die Frage, was die physikalischen
Konstanten sind und warum.

% ------------------------------------------------------------------------------
\section{Welches Framework ist vollständiger?}
% ------------------------------------------------------------------------------

Die Frage nach Vollständigkeit hängt davon ab, was ``vollständig'' bedeutet.

Wenn Vollständigkeit \emph{numerische Vorhersagekraft} bedeutet — die
Fähigkeit, spezifische Werte für Observablen aus minimaler Eingabe
abzuleiten — dann ist FFGFT vollständiger. Es erzeugt quantitative
Vorhersagen, die RSG ohne externe Parameterspezifikation nicht machen kann.

Wenn Vollständigkeit \emph{Allgemeinheit des Selektionsprinzips} bedeutet
— die Fähigkeit, zu charakterisieren welche Strukturen über verschiedene
zugrundeliegende Dynamiken hinweg persistieren — dann ist RSG breiter.
Sein Überlebensfilter gilt unabhängig von der spezifischen Feldtheorie;
FFGFTs Wicklungsstruktur ist eine besondere Realisierung einer persistenten
Konfiguration, keine allgemeine Beschreibung von Persistenz.

Die beiden Frameworks sind daher asymmetrisch komplementär:

\begin{itemize}
  \item FFGFT liefert die \emph{Skalenstruktur}, die RSGs Filter als Eingabe
        benötigt aber nicht ableitet.
  \item RSG liefert die \emph{Selektionsprinzip-Sprache}, die FFGFT
        instantiiert aber nicht als allgemeinen Operator formuliert.
\end{itemize}

Keines ist ein Overlay oder eine Brückenschicht des anderen. Sie sind
zwei eigenständige Frameworks, die angrenzende aber nicht identische
Fragen adressieren, mit partieller struktureller Überlappung auf der
Ebene der Persistenzbedingungen.

% ------------------------------------------------------------------------------
\section{Die Asymmetrie in einem Satz}
% ------------------------------------------------------------------------------

\begin{keyresult}[Strukturelle Asymmetrie]
RSG fragt: \emph{welche Histories überleben?} — ohne die Skala zu spezifizieren.
FFGFT fragt: \emph{was sind die Skalen?} — ohne einen allgemeinen
Persistenzoperator zu formulieren. RSG braucht FFGFTs Skalen um quantitativ
zu sein; FFGFT braucht RSGs Sprache um zu erklären warum Wicklungskonfigurationen
selektiert werden. Keines reduziert sich auf das andere.
\end{keyresult}

% ------------------------------------------------------------------------------
\section{Zur Visualisierung: zwei parallele Einträge}
% ------------------------------------------------------------------------------

Die strukturelle Analyse dieses Dokuments ergibt eine klare Anforderung
an jede vergleichende Visualisierung der beiden Frameworks.

FFGFT ist kein rechnerisches Overlay von RSG und keine Brückenschicht
innerhalb von RSGs Architektur. Es ist ein eigenständiges Framework mit
eigenem RA4-Primitiv ($T^4$-Torus), eigenem Lagrangian
($\mathcal{L}_{\text{T0}}$), eigenem Parameter ($\xi$) und eigenen
unabhängigen numerischen Vorhersagen, die RSG nicht liefert.

Eine strukturell korrekte Visualisierung zeigt daher:

\begin{itemize}
  \item \textbf{Zwei parallele Einträge auf derselben Strukturebene} —
        RSG und FFGFT als eigenständige Frameworks nebeneinander, nicht
        in einer Hierarchie.
  \item \textbf{Einen $\xi$-Pfad als eigene Spur} — der $\xi$-Parameter
        und die daraus abgeleiteten Skalen bilden eine eigenständige
        Linie in der Darstellung, getrennt von RSGs Filterparametern
        $\Gamma_i$ und $W_i$.
  \item \textbf{Korrespondenznotizen an den Berührungspunkten} — wo
        FFGFTs $\xi \cdot m_i$ den Wert von RSGs $\Gamma_i$ bestimmt,
        und wo RSGs Persistenzsprache das benennt, was FFGFTs
        Wicklungsstruktur implizit selektiert.
\end{itemize}

Diese Architektur — zwei gleichrangige Einträge mit expliziten
Korrespondenzpfeilen — ist die einzige Darstellung, die beiden Frameworks
gerecht wird: RSGs Allgemeinheit des Selektionsprinzips und FFGFTs
Vollständigkeit der Skalenableitung bleiben beide sichtbar, ohne dass
eines dem anderen untergeordnet wird.

% ------------------------------------------------------------------------------
\section{Zeichenerklärung}
% ------------------------------------------------------------------------------

{\small
\begin{tabular}{p{3cm}p{10.5cm}}
\toprule
\textbf{Symbol} & \textbf{Bedeutung} \\
\midrule
$S_i$ & RSG Überlebensgewicht von History $i$ \\
$\Gamma_i$ & RSG Verlustrate für History $i$ \\
$W_i$ & RSG Überlebensfaktor \\
$\xi = 4/30000$ & FFGFT dimensionsloser Kopplungsparameter; fixiert durch $\kappa=7$ \\
$T^4$ & 4-dimensionaler Torus; ontologisches Primitiv von FFGFT \\
$\kappa = 7$ & Eindeutiger Wicklungsexponent aus e-p-$\mu$-Massendreieck \\
$\mathcal{L}_{\text{T0}}$ & T0-Lagrangian (5 Terme; Dok.~180) \\
$\Delta m(x,t)$ & Dynamisches Massenfeld in T0 \\
$D_f = 3-\xi$ & Fraktale Dimension des physikalischen Raums ($\approx 2{,}99987$) \\
$L^2(T^4)$ & FFGFT Hilbertraum \\
$\alpha$ & Feinstrukturkonstante; aus $\xi$ hergeleitet (Dok.~009) \\
$\Lambda$ & Kosmologische Konstante aus $\xi$: $1{,}34\times10^{-122}$ (Dok.~182) \\
$H_0$ & Hubble-Konstante aus $\xi$: $66{,}2\,\text{km/s/Mpc}$ (Dok.~182) \\
RA4--RA1 & Schichten der Darstellungsarchitektur 2.1\\
          & \quad (Guevara Calderón 2026) \\
RSG & Recursive Survival Geometry (Austin 2024--2026) \\
FFGFT & Fundamentale Fraktale Geometrische Feldtheorie (Pascher 2024--2026) \\
\bottomrule
\end{tabular}
}

\vspace{1em}
\noindent\textbf{Schlüsseldokumente:}
Dok.~009 ($\xi$, $\kappa=7$) ·
Dok.~025 (CMB) ·
Dok.~026 (Rotverschiebung, FEM) ·
Dok.~180 (Lagrangian, $\kappa=7$) ·
Dok.~182 ($\Lambda$, $H_0$) ·
Dok.~230/231 (Hilbertraum-Brücke) ·
Dok.~245 (FFGFT RA-Diagnose) ·
\href{https://github.com/jpascher/T0-Time-Mass-Duality}{GitHub} ·
\href{https://doi.org/10.5281/zenodo.20117635}{Zenodo v1.1.0}
